\chapter{Practical 6: Flip-Flop Warriors: NAND Edition}

\section{Aim / Objective}
To design, construct, and validate the logic and operation of a JK Flip-Flop using discrete NAND gates (Universal Logic Gates).

\section{Apparatus / Components Required}

\begin{table}[htbp]
\centering
\small
\begin{tabularx}{\textwidth}{l X c}
\toprule
\textbf{Component} & \textbf{Specification / Type} & \textbf{Quantity} \\
\midrule
IC 7400 & Quad 2-Input NAND Gate & 1 \\
IC 7410 & Triple 3-Input NAND Gate & 1 \\
Breadboard & Standard & 1 \\
DC Power Supply & $+5\text{V}$ Regulated & 1 \\
LEDs & Red (Output $Q$), Green (Output $\bar{Q}$) & 2 \\
Resistors & $220\,\Omega$ or $330\,\Omega$ & 2 \\
Connecting Wires & Single strand / jumper wires & As req. \\
Logic Inputs & Toggle switches or DIP switches & 3 ($J, K, Clk$) \\
\bottomrule
\end{tabularx}
\caption{Apparatus and Components for Discrete JK Flip-Flop.}
\label{tab:ff-apparatus}
\end{table}

\section{Theory}

\subsection{The JK Flip-Flop}
The JK Flip-Flop is considered the "universal" flip-flop because it eliminates the invalid state found in SR flip-flops. It has two data inputs ($J$ and $K$) and a Clock input ($Clk$).
\begin{itemize}
    \item \textbf{$J$ (Jack):} Acts like the Set input.
    \item \textbf{$K$ (King):} Acts like the Reset input.
    \item \textbf{Feedback:} Outputs $Q$ and $\bar{Q}$ are fed back to inputs to enable toggling action.
\end{itemize}

\subsection{Logic Construction with NAND Gates}
To construct a JK Flip-Flop using discrete NAND gates, we require:
\begin{enumerate}
    \item \textbf{A Memory Cell (Latch):} Formed by two cross-coupled 2-input NAND gates (using IC 7400).
    \item \textbf{Steering / Input Logic:} Formed by two 3-input NAND gates (using IC 7410).
\end{enumerate}

\subsection{The Operation States}
\begin{enumerate}
    \item \textbf{Hold State ($J = 0, K = 0$):} No change in output ($Q_{n+1} = Q_n$).
    \item \textbf{Reset State ($J = 0, K = 1$):} Flip-flop resets ($Q = 0, \bar{Q} = 1$).
    \item \textbf{Set State ($J = 1, K = 0$):} Flip-flop sets ($Q = 1, \bar{Q} = 0$).
    \item \textbf{Toggle State ($J = 1, K = 1$):} Output switches to its complement ($Q_{n+1} = \bar{Q}_n$).
\end{enumerate}

\begin{commonmistake}[Race-Around Condition in Level-Triggered Setup]
In a basic level-triggered circuit constructed on a breadboard, the Toggle state ($J=1, K=1$) may result in a \textbf{Race-Around Condition} where the output oscillates rapidly while the clock is high. This confirms logic, though practical applications use Master-Slave configurations.
\end{commonmistake}

\section{Pin Diagrams and Connections}

\subsection{IC 7410 (Triple 3-Input NAND Gate)}
\begin{itemize}
    \item \textbf{Gate 1:} Inputs (Pins 1, 2, 13), Output (Pin 12)
    \item \textbf{Gate 2:} Inputs (Pins 3, 4, 5), Output (Pin 6)
    \item \textbf{Pin 7:} Ground ($\text{GND}$), \textbf{Pin 14:} Supply ($+5\text{V Vcc}$)
\end{itemize}

\section{Circuit Diagram / Setup}

\begin{figure}[htbp]
\centering
\begin{circuitikz}[scale=0.9, transform shape]
    % Steering NAND 1 (3-input)
    \node (nand3_1) at (0,2) [nand gate US, draw, logic gate inputs=nnn] {};
    % Steering NAND 2 (3-input)
    \node (nand3_2) at (0,-2) [nand gate US, draw, logic gate inputs=nnn] {};
    % Latch NAND 3 (2-input)
    \node (nand2_1) at (3.5,1.5) [nand gate US, draw, logic gate inputs=nn] {};
    % Latch NAND 4 (2-input)
    \node (nand2_2) at (3.5,-1.5) [nand gate US, draw, logic gate inputs=nn] {};
    
    % Connections
    \draw (nand3_1.output) -- (nand2_1.input 1);
    \draw (nand3_2.output) -- (nand2_2.input 2);
    
    \draw (nand2_1.output) -- ++(0.5,0) coordinate (q_out) -- ++(1,0) node[right] {$Q$};
    \draw (nand2_2.output) -- ++(0.5,0) coordinate (qbar_out) -- ++(1,0) node[right] {$\bar{Q}$};
    
    % Cross-coupling
    \draw (q_out) -- ++(0,-0.8) -- ($(nand2_2.input 1)+(-0.3,0)$) -- (nand2_2.input 1);
    \draw (qbar_out) -- ++(0,0.8) -- ($(nand2_1.input 2)+(-0.3,0)$) -- (nand2_1.input 2);
    
    % Inputs
    \draw (nand3_1.input 1) -- ++(-1,0) node[left] {$J$};
    \draw (nand3_1.input 2) -- ++(-1,0) node[left] {$Clk$};
    \draw (nand3_2.input 2) -- ++(-1,0) node[left] {$Clk$};
    \draw (nand3_2.input 3) -- ++(-1,0) node[left] {$K$};
    
    % Feedbacks
    \draw (qbar_out) -- ++(0,2.3) -- ($(nand3_1.input 3)+(-0.5,2.7)$) -- ($(nand3_1.input 3)+(-0.5,0)$) -- (nand3_1.input 3);
    \draw (q_out) -- ++(0,-2.3) -- ($(nand3_2.input 1)+(-0.5,-2.7)$) -- ($(nand3_2.input 1)+(-0.5,0)$) -- (nand3_2.input 1);
\end{circuitikz}
\caption{Discrete JK Flip-Flop Circuit Schematic using 2-Input and 3-Input NAND Gates.}
\label{fig:jk-flipflop-schematic}
\end{figure}

\section{Step-by-Step Procedure}

\begin{enumerate}
    \item \textbf{Placement:} Insert IC 7400 and IC 7410 into breadboard groove firmly.
    \item \textbf{Power:} Connect Pin 14 of both ICs to $+5\text{V (Vcc)}$ and Pin 7 of both ICs to Ground ($\text{GND}$).
    \item \textbf{Latch Construction (Using IC 7400):}
    \begin{itemize}
        \item Connect Output Pin 3 to Input Pin 5.
        \item Connect Output Pin 6 to Input Pin 2.
        \item Connect Pin 1 to output of J-steering gate.
        \item Connect Pin 4 to output of K-steering gate.
    \end{itemize}
    \item \textbf{Steering Logic (Using IC 7410):}
    \begin{itemize}
        \item J-Gate: Use pins 1, 2, 13 (Inputs) and 12 (Output). Connect Pin 12 to Pin 1 of 7400.
        \item K-Gate: Use pins 3, 4, 5 (Inputs) and 6 (Output). Connect Pin 6 to Pin 4 of 7400.
    \end{itemize}
    \item \textbf{Feedback Connections:}
    \begin{itemize}
        \item Connect Output $Q$ (Pin 3 of 7400) back to input of K-Gate (Pin 5 of 7410).
        \item Connect Output $\bar{Q}$ (Pin 6 of 7400) back to input of J-Gate (Pin 13 of 7410).
    \end{itemize}
    \item \textbf{Inputs:} Connect switches to J (Pin 1 of 7410), K (Pin 3 of 7410), and Clock (Pin 2 and Pin 4 of 7410 tied together).
    \item \textbf{Outputs:} Connect LED (via $330\,\Omega$ resistor) to Pin 3 of 7400 ($Q$) and Pin 6 of 7400 ($\bar{Q}$).
    \item \textbf{Validation:} Apply logic combinations as per observation table and record LED status.
\end{enumerate}

\section{Observations and Truth Table}

\noindent \textbf{Notation:} $Q_n$: Current State $\quad \bullet \quad$ $Q_{n+1}$: Next State $\quad \bullet \quad$ Logic $0$ = OFF ($0\text{V}$), Logic $1$ = ON ($+5\text{V}$)

\begin{table}[htbp]
\centering
\small
\begin{tabularx}{\textwidth}{c c c c c X}
\toprule
\textbf{CLOCK ($CLK$)} & \textbf{INPUT $J$} & \textbf{INPUT $K$} & \textbf{PRESENT ($Q_N$)} & \textbf{NEXT ($Q_{N+1}$)} & \textbf{ACTION / MODE} \\
\midrule
0 & X & X & 0 & 0 & No Change \\
0 & X & X & 1 & 1 & No Change \\
$\uparrow (1)$ & 0 & 0 & 0 & 0 & \textbf{Hold} \\
$\uparrow (1)$ & 0 & 0 & 1 & 1 & \textbf{Hold} \\
$\uparrow (1)$ & 0 & 1 & 0 & 0 & \textbf{Reset} \\
$\uparrow (1)$ & 0 & 1 & 1 & 0 & \textbf{Reset} \\
$\uparrow (1)$ & 1 & 0 & 0 & 1 & \textbf{Set} \\
$\uparrow (1)$ & 1 & 0 & 1 & 1 & \textbf{Set} \\
$\uparrow (1)$ & 1 & 1 & 0 & 1 & \textbf{Toggle} \\
$\uparrow (1)$ & 1 & 1 & 1 & 0 & \textbf{Toggle} \\
\bottomrule
\end{tabularx}
\caption{JK Flip-Flop State Transition Truth Table.}
\label{tab:jk-truth-table}
\end{table}

\section{Result}
The JK Flip-Flop was successfully designed using NAND gates (IC 7400 and IC 7410). Logic states (Hold, Set, Reset, Toggle) were verified against the truth table.

\textbf{Conclusion on Toggle State:} When $J=1$ and $K=1$, output toggles. However, due to continuous clock pulse in this discrete setup, a "Race-Around" condition may be observed (LEDs appearing dim or flickering) indicating output is oscillating while clock is high.

\section{Viva Questions \& Answers}

\begin{vivaqa}[JK Flip-Flop Universality]
\textbf{Question:} Why is the JK Flip-Flop called a "Universal" Flip-Flop?\\[4pt]
\textbf{Answer:} It is versatile and solves the indeterminate/invalid state problem of the SR flip-flop. It can also be configured to act as a D flip-flop or T flip-flop.
\end{vivaqa}

\begin{vivaqa}[Race-Around Condition Definition]
\textbf{Question:} What is the "Race-Around Condition"?\\[4pt]
\textbf{Answer:} In a level-triggered JK flip-flop, if $J=1, K=1$ and the clock pulse width ($t_p$) is greater than gate propagation delay, the output toggles multiple times within a single clock pulse (unwanted oscillation).
\end{vivaqa}

\begin{vivaqa}[Eliminating Race-Around Condition]
\textbf{Question:} How do we eliminate the Race-Around Condition?\\[4pt]
\textbf{Answer:} By using a \textbf{Master-Slave JK Flip-Flop} configuration or by using edge-triggered flip-flops where pulse width is effectively negligible.
\end{vivaqa}

\begin{vivaqa}[Preset and Clear Pins]
\textbf{Question:} What is the function of Clear ($\text{CLR}$) and Preset ($\text{PRE}$) pins found on standard ICs (like 7476)?\\[4pt]
\textbf{Answer:} These are asynchronous inputs. $\text{PRE}$ sets $Q=1$ immediately irrespective of Clock; $\text{CLR}$ resets $Q=0$ immediately irrespective of Clock.
\end{vivaqa}

\begin{vivaqa}[2-Input NAND Realization]
\textbf{Question:} Can we realize this circuit using only IC 7400 (2-input NANDs)?\\[4pt]
\textbf{Answer:} Yes, but it is more complex. You would need to construct 3-input NAND functions by cascading two 2-input NAND gates, requiring more gates and wiring.
\end{vivaqa}
